L'acide sulfurique à l'état pur est un liquide visqueux et très soluble dans
l'eau. Sa solution est utilisée comme électrolyte dans les accumulateurs au
plomb constituant les batteries automobiles. La concentration molaire de la
solution est alors de $\qty{6.0}{\mol\per\L}$. L'équation de la réaction de
dissolution de l'acide sulfurique dans l'eau s'écrit~:

\begin{center}
    \ce{H2SO4$(\ell)$ +  2H2O -> 2H3O+ + SO$_4^{2-}$(aq)}
\end{center}

\begin{enumerate}
    \item Calculer les concentrations $\left[\ce{H3O+}\right]$ et
          $\left[\ce{SO$_4^{2-}$(aq)}\right]$.
    \item Comment obtenir à partir de la solution précédente $\qty{500}{\mL}$ de
          solution telle que dans cette solution
          $\left[\ce{SO$_4^{2-}$(aq)}\right]=\qty{1.0}{\mol\per\L}$~?
    \item Comment obtenir, à partir de la solution obtenue précédemment,
          $\qty{500}{\mL}$ de solution telle que, dans celle-ci~:
          $\left[\ce{H3O+}\right]=\qty{0.02}{\mol\per\L}$~?
\end{enumerate}

